\begin{abstract}
    Consider the following variation on the classical \textsc{Hierarchical Clustering} problem: Usually, while building a hierarchical clustering, one recursively partitions the data until each cluster becomes a singleton. In our setup, we relax the halting condition of the recursive process to stop whenever the remaining cluster is a graph belonging to a given class $\mathcal{F}$. We call this problem \textsc{Hierarchical $\mathcal{F}$-Clustering} and we measure the quality of any solution using adapted Dasgupta's clustering objective. We study two natural choices of $\mathcal{F}$: trees, denoted by $\mathcal{T}$ and graphs of diameter bounded by~$d$, denoted by $\mathcal{D}_d$ (which includes cliques for $d=1$).

    Hereby, we present the first polynomial time $\mathcal{O}(\log n\cdot\log\log n)$ and $\mathcal{O}(\log n)$-approximation algorithms for clustering into $\mathcal{T}$- and $\mathcal{D}_d$-clusters respectively.
    Our main technical contribution is a framework for approximating such problems based on linear programming. In fact, we characterize graphs classes $\mathcal{F}$ for which our approach can be applied and show that it includes both trees and bounded diameter graphs. However, our ideas are not limited to them and might be useful for other structures as well. Broadly speaking, our framework applies whenever the corresponding flat clustering problem, which we call \textsc{$p_{\mathcal{F}}$-Partitioning}, admits a natural ILP formulation together with a rounding procedure with provable approximation guarantees. Intuitively, given a set of vertices called terminals, the problem is to find an edge set whose removal results in satisfying certain vertex-dependent structural predicate for each terminal. For example, excluding all cycles going through terminals or ensuring that the diameter of all clusters with respect to the terminals is not too large. We then use these ingredients to build clustering trees with the aforementioned approximation guarantees. To complement these results, we show that both \textsc{Hierarchical $\mathcal{T}$-Clustering} and \textsc{Hierarchical $\mathcal{D}_d$-Clustering} cannot be approximated within any constant factor under the \textsc{Small Set Expansion Hypothesis}.
\end{abstract}
